\documentclass[11pt]{amsart}
\usepackage[a4paper,margin=1in]{geometry}
\usepackage{amsmath,amssymb,amsthm}
\usepackage{booktabs}
\usepackage{pgfplots}
\pgfplotsset{compat=1.18}
\usepackage[hidelinks]{hyperref}
\newcommand{\doi}[1]{doi:\href{https://doi.org/#1}{#1}}

\newtheorem{thm}{Theorem}[section]
\newtheorem{lemma}[thm]{Lemma}
\newtheorem{prop}[thm]{Proposition}
\newtheorem{cor}[thm]{Corollary}
\newtheorem{conj}[thm]{Conjecture}
\theoremstyle{definition}
\newtheorem{definition}[thm]{Definition}
\newtheorem{remark}[thm]{Remark}

\DeclareMathOperator{\corr}{corr}
\DeclareMathOperator{\cov}{cov}
\DeclareMathOperator{\var}{var}
\newcommand{\E}{\mathbb{E}}
\newcommand{\rs}{\rho_S}

\title[Discretising a Frank copula can inflate correlation]{Discretising a Frank copula can inflate linear correlation}
\author{Vamshi Jandhyala}
\address{London, United Kingdom}
\date{Preprint, version 1, 24 November 2026. Not peer reviewed.}
\subjclass[2020]{62H20, 60E15, 62H05}
\keywords{Frank copula, ordinal data, correlation attenuation, discretisation, monotone correlation}

\begin{document}

\begin{abstract}
Cut a pair of uniform random variables joined by a Frank copula into ordinal categories, score the categories
$1,2,\dots,m$, and compare the Pearson correlation of the scores with Spearman's rho of the latent pair. Barbiero
conjectured, from extensive numerical evidence with at most eleven categories, that the scores are never more
correlated than the latent pair. We show that the conjecture fails for every nonzero value of the Frank parameter.
The reason is that the regression of one latent variable on the other is never affine, so a slight monotone
distortion of the scale correlates better than the scale itself, and fine enough cut-points reproduce the gain. An
explicit example with fifteen categories and latent Spearman's rho $0.8602$ gives implied correlation $0.8604$,
certified by interval arithmetic. The key steps are formalised in Lean.
\end{abstract}

\maketitle

\section{Introduction}

A common model for two ordinal survey answers takes a latent pair of continuous variables and cuts each one at fixed
thresholds. The Pearson correlation of the resulting integer scores, the \emph{implied} correlation, is then compared
with the dependence of the latent pair. For a bivariate normal latent pair this comparison has a classical answer:
cutting never raises the absolute correlation, because the normal distribution's maximal correlation equals its
correlation coefficient~\cite{Lancaster}. Barbiero~\cite{Barbiero} proposed the Frank copula as an alternative latent
model and asked whether it shares this property, measuring latent dependence by Spearman's rho $\rs$, the
correlation of the uniform latent variables themselves.

Here is the setting in full. Let $(U,V)$ have the Frank copula $C_\theta$ with parameter $\theta\ne0$,
\begin{equation}\label{eq:frank}
C_\theta(u,v)=-\frac1\theta\log\Bigl(1+\frac{(e^{-\theta u}-1)(e^{-\theta v}-1)}{e^{-\theta}-1}\Bigr),
\qquad 0\le u,v\le1,
\end{equation}
so that $U$ and $V$ are uniform on $[0,1]$ and $\corr(U,V)=\rs(\theta)$. Choose cut-points $0<a_1<\dots<a_{m-1}<1$
and $0<b_1<\dots<b_{n-1}<1$ and put
\[
X=1+\#\{k : U>a_k\},\qquad Y=1+\#\{l : V>b_l\}.
\]
Then $X$ is an ordinal variable on $\{1,\dots,m\}$ with $P(X\le j)=a_j$, and every law on $\{1,\dots,m\}$
that gives each category positive probability arises this way. Barbiero showed that $\corr(X,Y)$ has the sign of $\rs(\theta)$~\cite[Theorem~2]{Barbiero},
and on the basis of a numerical study over $m\le11$ categories and $\rs\le0.9$ conjectured more.

\begin{conj}[Barbiero~{\cite[Sections~2.2 and~5]{Barbiero}}]\label{conj}
Discretisation of a Frank copula latent variable cannot inflate linear correlation: $|\corr(X,Y)|\le|\rs(\theta)|$
for all cut-points, with equality only in limiting cases.
\end{conj}

Barbiero also reports, from joint work with Hitaj, that a Student~$t$ latent pair can inflate correlation after
discretisation when the degrees of freedom and the latent correlation are small and the margins are strongly
skewed~\cite[Section~2.2]{Barbiero}. Loginov~\cite{Loginov} studies the same cutting scheme with the copula fixed and
the thresholds free, bounding the error of treating the grades as independent; that is a statement about means of
functions of the grades, and it does not address the conjecture. We show that the conjecture is false for every
$\theta\ne0$.

\begin{thm}\label{thm:main}
For every $\theta\ne0$ there are cut-points for which $|\corr(X,Y)|>|\rs(\theta)|$.
\end{thm}

A concrete instance comes first. Take $\theta=10$, so $\rs=0.86023$, and fifteen categories on both sides, with the
same cut-points
\[
a=b=\Bigl(\tfrac{27}{250},\tfrac{193}{1000},\tfrac{13}{50},\tfrac{8}{25},\tfrac{187}{500},\tfrac{17}{40},
\tfrac{19}{40},\tfrac{21}{40},\tfrac{23}{40},\tfrac{313}{500},\tfrac{17}{25},\tfrac{37}{50},\tfrac{807}{1000},
\tfrac{223}{250}\Bigr).
\]
The category probabilities are symmetric and U-shaped (Table~\ref{tab:example}), the shape Barbiero's optimiser
found for small $m$. The implied correlation is $0.86039$, which exceeds $\rs$ by $1.53\times10^{-4}$
(Proposition~\ref{prop:cert}).

\begin{table}[h]
\centering\footnotesize\setlength{\tabcolsep}{4pt}
\begin{tabular}{@{}l*{15}{c}@{}}
\toprule
$j$ & 1&2&3&4&5&6&7&8&9&10&11&12&13&14&15\\
\midrule
$P(X=j)$ & .108&.085&.067&.060&.054&.051&.050&.050&.050&.051&.054&.060&.067&.085&.108\\
\bottomrule
\end{tabular}
\caption{The ordinal law of $X$ (and of $Y$) in the fifteen-category counterexample, $\theta=10$.}\label{tab:example}
\end{table}

The gain is small and needs many categories, which explains why a study up to eleven categories saw only
attenuation. Cutting a continuous variable into $m$ pieces costs correlation roughly in proportion to $m^{-2}$ (an
empirical rate in our computations, not used in any proof), while the gain available from a better scale is fixed. Section~\ref{sec:numerics} records where the balance tips: at
$\rs=0.86$ the first inflating discretisation we found has fifteen categories, and at $\rs=0.5$ it has
forty-eight.

The mechanism is a property of the latent copula rather than of the cut-points. Write $r(u)=\E[V\mid U=u]$. If the
identity scale gave the largest correlation among increasing transforms of $U$, then $r$ would have to be affine
(Lemma~\ref{lem:var}). For the Frank copula $r$ has a closed form (Proposition~\ref{prop:reg}), and three of its
values are never collinear (Lemma~\ref{lem:notaffine}). Figure~\ref{fig:reg} shows how far $r$ is from its best affine
fit. Integer scores of a fine discretisation of the slightly distorted scale then carry the gain over to the
ordinal variables.

\begin{figure}[t]
\centering
\begin{tikzpicture}
\begin{axis}[width=0.7\textwidth,height=5.2cm,xlabel={$u$},ylabel={$r(u)-\bigl(\tfrac12+\rs(u-\tfrac12)\bigr)$},
  xmin=0,xmax=1,ymin=-0.016,ymax=0.016,grid=major,legend pos=north east,legend style={font=\small},
  yticklabel style={/pgf/number format/fixed,/pgf/number format/precision=3},scaled y ticks=false]
\addplot[thick,blue] table[x=u,y=gap] {fig_reg_t10.dat};
\addplot[thick,red,dashed] table[x=u,y=gap] {fig_reg_t3.dat};
\addplot[black,thin,forget plot] coordinates {(0,0) (1,0)};
\legend{$\theta=10$,$\theta=3$}
\end{axis}
\end{tikzpicture}
\caption{The regression $r(u)=\E[V\mid U=u]$ of the Frank copula minus its best affine approximation. The curve
is odd about $u=\tfrac12$ because the Frank copula is radially symmetric. A nonzero curve means some increasing
rescaling of $U$ correlates with $V$ better than $U$ does.}\label{fig:reg}
\end{figure}

\section{Moments of ordinal scores}

\begin{lemma}\label{lem:moments}
Let $(U,V)$ have uniform margins and copula $C$, and let $X,Y$ be the scores defined above. Then
\[
\cov(X,Y)=\sum_{k=1}^{m-1}\sum_{l=1}^{n-1}\bigl(C(a_k,b_l)-a_kb_l\bigr),\qquad
\var X=\sum_{k=1}^{m-1}\sum_{k'=1}^{m-1}\bigl(\min(a_k,a_{k'})-a_ka_{k'}\bigr).
\]
\end{lemma}

\begin{proof}
$X-1$ is a sum of the indicators $\mathbf 1\{U>a_k\}$, so the covariance is a double sum of indicator covariances.
By inclusion and exclusion $P(U>a,V>b)=1-a-b+C(a,b)$, and subtracting $P(U>a)P(V>b)=(1-a)(1-b)$ leaves $C(a,b)-ab$.
The variance is the case $V=U$, where $C(a,a')=P(U\le\min(a,a'))=\min(a,a')$.
\end{proof}

Read in words, each pair of cut-points contributes the amount by which the copula exceeds independence at that
point, and the variance is the same sum for the comonotone copula $\min(u,v)$.

\section{The regression of the Frank copula}

Fix $\theta\ne0$ and write
\[
w=e^{-\theta u},\qquad x=w-1,\qquad a=e^{-\theta}-1.
\]
Differentiating~\eqref{eq:frank} in $u$ gives the conditional distribution function of $V$ given $U=u$,
\begin{equation}\label{eq:K}
K(u,v)=\frac{\partial C_\theta}{\partial u}(u,v)=\frac{w\,(e^{-\theta v}-1)}{a+x\,(e^{-\theta v}-1)}.
\end{equation}
The denominator does not vanish for $0\le u,v\le1$: it equals $-\bigl((e^{-\theta v}-e^{-\theta})+w(1-e^{-\theta v})\bigr)$,
and the two terms in the bracket have the sign of $\theta$ and do not both vanish. So $K$ is smooth on a
neighbourhood of the unit square, and $r(u)=\E[V\mid U=u]=1-\int_0^1K(u,v)\,dv$ is smooth on $[0,1]$.

\begin{prop}\label{prop:reg}
For $0<u<1$,
\[
r(u)=1-\frac{w\,(x-au)}{x\,(x-a)}.
\]
\end{prop}

We make no claim that this formula is new: it follows from~\eqref{eq:K} by direct integration, and we include the
derivation because the rest of the argument rests on it.

\begin{proof}
For fixed $u$ the function
\[
G(v)=\frac{w}{x}\Bigl(v-\frac{a}{a-x}\Bigl(v+\frac1\theta\log\bigl|a+x(e^{-\theta v}-1)\bigr|\Bigr)\Bigr)
\]
has $G'(v)=K(u,v)$: the bracket has derivative $(a-x)/(a+x(e^{-\theta v}-1))$, so
$G'=\frac wx\bigl(1-\frac{a}{a+x(e^{-\theta v}-1)}\bigr)$, which is~\eqref{eq:K}. Here $x\ne0$ because $u>0$, and
$a\ne x$ because $u<1$. The argument of the logarithm is $a$ at $v=0$ and $wa$ at $v=1$, so the logarithm increases by
$\log w=-\theta u$, and
\[
\int_0^1K(u,v)\,dv=G(1)-G(0)=\frac wx\Bigl(1-\frac{a(1-u)}{a-x}\Bigr)=\frac{w\,(x-au)}{x\,(x-a)}.\qedhere
\]
\end{proof}

The numerator $x-au$ compares the value $e^{-\theta u}-1$ with its chord from $u=0$ to $u=1$, and the factor
$x(x-a)$ vanishes exactly at those two endpoints; both singularities are removable, with $r(0)$ and $r(1)$ given by
the limits. Nothing in the formula is linear in $u$ except the chord, which is the first hint that $r$ is not affine.

\begin{lemma}\label{lem:notaffine}
For every $\theta\ne0$ the function $r$ is not affine on $(0,1)$.
\end{lemma}

\begin{proof}
An affine function $\ell$ satisfies $\ell(\tfrac12)-\ell(\tfrac14)=2\bigl(\ell(\tfrac14)-\ell(\tfrac18)\bigr)$.
Put $s=e^{-\theta/8}$, so that $w=s^k$ at $u=k/8$ and $e^{-\theta}=s^8$. Substituting into
Proposition~\ref{prop:reg} and clearing denominators gives
\[
\bigl(r(\tfrac12)-r(\tfrac14)\bigr)-2\bigl(r(\tfrac14)-r(\tfrac18)\bigr)
=\frac{-s\,(s-1)^3(s+1)(s^4+s^3+3s^2+s+1)}{4\,(s^2-s+1)(s^2+s+1)(s^6+s^5+s^4+s^3+s^2+s+1)}.
\]
For $s>0$ every factor except $(s-1)^3$ is positive, and $s\ne1$ because $\theta\ne0$.
\end{proof}

Two other proofs are worth knowing. The closed form is meromorphic in $u$ with a simple pole at $u=2\pi i/\theta$,
which an affine function cannot have. Alternatively, differentiating under the integral twice gives
\[
r''(0)=\frac{\theta}{2(1-e^{-\theta})^3}\int_{e^{-\theta}}^1\frac{(1-y)^2(y-e^{-\theta})^2}{y^2}\,dy,
\]
which has the sign of $\theta$. The three-point argument is the one formalised, because it is pure algebra.

\section{From a better scale to better categories}

\begin{lemma}\label{lem:var}
Let $\xi,\eta,\zeta$ be vectors in a real inner product space with $\langle\xi,\eta\rangle>0$,
$\langle\zeta,\xi\rangle=0$ and $\langle\zeta,\eta\rangle>0$. Then there is $\varepsilon_0>0$ such that
$\corr(\xi+\varepsilon\zeta,\eta)>\corr(\xi,\eta)$ for every $0<\varepsilon\le\varepsilon_0$, where
$\corr(\xi,\eta)=\langle\xi,\eta\rangle/(\|\xi\|\,\|\eta\|)$.
\end{lemma}

\begin{proof}
Write $c=\langle\xi,\eta\rangle$, $d=\langle\zeta,\eta\rangle$, $n=\|\xi\|$ and $q=\|\zeta\|$. Orthogonality gives
$\|\xi+\varepsilon\zeta\|^2=n^2+\varepsilon^2q^2$, so the claim is $(c+\varepsilon d)^2n^2>c^2(n^2+\varepsilon^2q^2)$,
that is $\varepsilon\bigl(2cdn^2+\varepsilon(d^2n^2-c^2q^2)\bigr)>0$. This holds for
$0<\varepsilon\le\varepsilon_0=cdn^2/(c^2q^2+1)$, since then $\varepsilon c^2q^2<cdn^2$.
\end{proof}

\begin{proof}[Proof of Theorem~\ref{thm:main}]
First let $\theta>0$, so $\rho=\rs(\theta)>0$. Since $U$ and $V$ are uniform, the best affine approximation to $r$ in
$L^2(U)$ is $p(u)=\tfrac12+\rho(u-\tfrac12)$: indeed $\E\,r(U)=\E V=\tfrac12$ and
$\cov(U,r(U))=\cov(U,V)=\rho/12$. Let $h=r-p$. Then $\E\,h(U)=0$ and $\E\bigl[(U-\tfrac12)h(U)\bigr]=0$, and by the
tower property
\[
\E\bigl[h(U)(V-\tfrac12)\bigr]=\E\bigl[h(U)\bigl(h(U)+\rho(U-\tfrac12)\bigr)\bigr]=\E\,h(U)^2>0,
\]
the last because $h$ is continuous and, by Lemma~\ref{lem:notaffine}, not identically zero on $(0,1)$. Apply
Lemma~\ref{lem:var} in $L^2$ with $\xi=U-\tfrac12$, $\eta=V-\tfrac12$ and $\zeta=h(U)$. Since $h'$ is bounded on
$[0,1]$, we may also take $\varepsilon\sup|h'|<1$, so that $f(u)=u+\varepsilon h(u)$ is continuous and strictly
increasing and $\corr(f(U),V)>\rho$.

Now discretise. For a positive integer $M$ put $X_M=\lfloor Mf(U)\rfloor$ and $Y_M=\lfloor MV\rfloor$. Then
$|X_M/M-f(U)|\le1/M$ and $|Y_M/M-V|\le1/M$, so $\corr(X_M,Y_M)\to\corr(f(U),V)$ as $M\to\infty$, the correlation being
continuous in $L^2$ away from constant variables. Fix $M$ with $\corr(X_M,Y_M)>\rho$. Because $f$ is continuous and
increasing, $X_M$ is a nondecreasing function of $U$ whose values of positive probability form a run of
consecutive integers; the only other possible value is $\lfloor Mf(1)\rfloor$ at the null event $U=1$, which we
discard. After subtracting the least value of positive probability, $X_M$ is therefore almost surely an ordinal
variable of the form $X$ above, with $a_j=P(X_M\le x_0+j-1)$ for that least value $x_0$; the same holds for $Y_M$,
whose value $M$ at $V=1$ is likewise null. Shifts do not change correlation, which proves the theorem
for $\theta>0$.

For $\theta<0$ use the reflection $C_\theta(u,v)=u-C_{-\theta}(u,1-v)$: if $(U,V)$ has copula $C_\theta$ then
$(U,1-V)$ has copula $C_{-\theta}$, and $\rs(\theta)=-\rs(-\theta)$. Take the inflating scores $X,Y'$ for $(U,1-V)$
with $n$ categories and set $Y=n+1-Y'$, a nondecreasing integer score of $V$. Then
$|\corr(X,Y)|=\corr(X,Y')>\rs(-\theta)=|\rs(\theta)|$.
\end{proof}

\begin{cor}
For every $\theta\ne0$ the maximal monotone correlation of the Frank copula~\cite{KS}, the supremum of
$|\corr(f(U),g(V))|$ over monotone $f,g$, strictly exceeds $|\rs(\theta)|$.
\end{cor}

\begin{remark}
The argument uses nothing about the Frank copula beyond the smoothness of $r$ and Lemma~\ref{lem:notaffine}. Any
copula whose regression function is smooth and not affine admits an inflating discretisation, measured against its
Spearman's rho. An attenuation property of the kind conjectured therefore needs affine regressions in both
directions. The Farlie--Gumbel--Morgenstern copula, with density $1+\alpha(1-2u)(1-2v)$, has them, and there the property holds:
for square-integrable $f,g$, $\cov(f(U),g(V))=\alpha\,\langle f,\phi\rangle\langle g,\phi\rangle$ with
$\phi(u)=1-2u$ and $\|\phi\|^2=\tfrac13$, so the Cauchy--Schwarz inequality gives
$|\corr(f(U),g(V))|\le|\alpha|/3=|\rs|$. The normal copula's regression of $V$ on $U$ is not affine either, and indeed the normal scores
$\Phi^{-1}(U),\Phi^{-1}(V)$ have correlation $\rho>\rs=\tfrac6\pi\arcsin\tfrac\rho2$ for $0<\rho<1$; the classical
attenuation result is measured against $\rho$, not against $\rs$. For the Frank copula the sharp bound over all
discretisations is the maximal monotone correlation; whether it has a closed form is open.
\end{remark}

\section{The certified example and the number of categories}\label{sec:numerics}

\begin{prop}\label{prop:cert}
For $\theta=10$ and the cut-points of the introduction,
\begin{align*}
\corr(X,Y)&\in0.8603867893143344074910842\pm9\times10^{-27},\\
\rs(10)&\in0.8602336388082110151918579\pm1.1\times10^{-26}.
\end{align*}
In particular $\corr(X,Y)-\rs(10)=1.531505061\ldots\times10^{-4}>0$.
\end{prop}

\begin{proof}
Lemma~\ref{lem:moments} reduces $\corr(X,Y)$ to $196$ values of~\eqref{eq:frank} at rational points and an exact
rational variance, $11727/500$. For $\rs$ we use Genest's formula~\cite{Genest}
$\rs(\theta)=1-\tfrac{12}{\theta}\bigl(D_1(\theta)-D_2(\theta)\bigr)$ with the Debye functions
$D_k(\theta)=\tfrac{k}{\theta^k}\int_0^\theta\tfrac{t^k}{e^t-1}\,dt$. Both quantities were evaluated in Arb ball
arithmetic~\cite{Arb} at 200 bits; the integral over $[2^{-100},\theta]$ is computed by rigorous quadrature and the
piece over $[0,2^{-100}]$ is bounded by $2^{-100k}/k$, since $0<t/(e^t-1)\le1$. An independent check with separate
code, using 70-digit interval arithmetic and the series $\int_0^\theta t^k/(e^t-1)\,dt$ in terms of $\zeta(2)$ and
$\zeta(3)$ with explicit tail bounds, gives the gap in $[1.53150506099\times10^{-4},1.53150506148\times10^{-4}]$.
\end{proof}

Table~\ref{tab:mstar} lists, for several values of $\rs$, the smallest number of categories at which a numerical
optimisation over cut-points (both variables with $m$ categories) found inflation. These values come from a local
optimiser with restarts and are upper bounds for the true minimum; only the example above is certified.
Figure~\ref{fig:gap} shows the largest implied correlation found against $m$ at $\theta=10$: no inflation was found
for $m\le14$, and from $m=15$ the gain grows towards the limit set by the maximal monotone correlation. A
floating-point computation of the largest correlation of increasing scores on an equiprobable $400\times400$
discretisation gives $0.8633$, numerical evidence that the limit is near that value; it is not certified.

\begin{table}[h]
\centering\small
\begin{tabular}{@{}lccccccc@{}}
\toprule
$\rs$ & 0.5 & 0.6 & 0.7 & 0.8 & 0.86 & 0.9 & 0.95\\
$\theta$ & 3.446 & 4.466 & 5.821 & 7.902 & 9.990 & 12.26 & 18.20\\
\midrule
least $m$ found & 48 & 32 & 23 & 17 & 15 & 14 & 14\\
\bottomrule
\end{tabular}
\caption{Fewest categories at which inflation was found numerically.}\label{tab:mstar}
\end{table}

\begin{figure}[t]
\centering
\begin{tikzpicture}
\begin{axis}[width=0.7\textwidth,height=5cm,xlabel={number of categories $m$},
  ylabel={largest $\corr(X,Y)-\rs$ found},xmin=2,xmax=40,ymin=-0.02,ymax=0.004,grid=major,
  yticklabel style={/pgf/number format/fixed,/pgf/number format/precision=3},scaled y ticks=false]
\addplot[thick,blue,mark=*,mark size=1pt] table[x=m,y=gap] {fig_gap_t10.dat};
\addplot[black,thin] coordinates {(2,0) (40,0)};
\end{axis}
\end{tikzpicture}
\caption{Largest implied correlation found by local optimisation, minus $\rs$, at $\theta=10$ ($\rs=0.8602$), over
cut-points with $m$ categories on each variable. Values for $m\le5$ lie below the plotted range. Inflation is first
found at $m=15$; for $m\le14$ none was found, which does not prove that none exists.}\label{fig:gap}
\end{figure}

\section{Data, code and the use of AI tools}\label{sec:data}
\sloppy

The scripts, the certificates and the Lean files are archived on Zenodo, \doi{10.5281/zenodo.24497009}. The same
files are in the folder \texttt{frank-copula-inflation} of the repository
\href{https://github.com/jvvk/mathematics}{github.com/jvvk/mathematics}, at tag \texttt{frank-copula-inflation-v1}.

Parts of the argument are formalised in Lean~4 with Mathlib: the derivative~\eqref{eq:K}, the integral of
Proposition~\ref{prop:reg}, Lemma~\ref{lem:notaffine}, Lemma~\ref{lem:var} and Lemma~\ref{lem:moments} for an
arbitrary copula. Not formalised are the identity $r(u)=1-\int_0^1K\,dv$, the tower-property step and the
approximation argument in the proof of Theorem~\ref{thm:main}, which are standard, and the numerical certificate,
which is checked by two independent programs.

The author used AI tools. Claude (Anthropic) found the counterexample, the pole argument and the three-point proof,
derived the closed form of Proposition~\ref{prop:reg}, and wrote the code, the Lean formalisation and the text. Codex (OpenAI) independently rechecked the
counterexample and the variational lemma, and found the proof through $r''(0)$. The author is responsible for the
content.

\begin{thebibliography}{9}
\bibitem{Barbiero} A.~Barbiero, The Frank copula as a latent model for ordinal variables: dependence structure,
correlation attenuation, and statistical inference, \emph{Adv.\ Data Anal.\ Classif.} (2026),
\doi{10.1007/s11634-026-00709-1}.
\bibitem{Genest} C.~Genest, Frank's family of bivariate distributions, \emph{Biometrika} \textbf{74} (1987),
549--555, \doi{10.1093/biomet/74.3.549}.
\bibitem{Arb} F.~Johansson, Arb: efficient arbitrary-precision midpoint-radius interval arithmetic, \emph{IEEE
Trans.\ Comput.} \textbf{66} (2017), 1281--1292, \doi{10.1109/TC.2017.2690633}.
\bibitem{KS} G.~Kimeldorf and A.~R.~Sampson, Monotone dependence, \emph{Ann.\ Statist.} \textbf{6} (1978),
895--903, \doi{10.1214/aos/1176344262}.
\bibitem{Lancaster} H.~O.~Lancaster, Some properties of the bivariate normal distribution considered in the form of
a contingency table, \emph{Biometrika} \textbf{44} (1957), 289--292, \doi{10.2307/2333274}.
\bibitem{Loginov} K.~V.~Loginov, Worst-case bias of the independence assumption for functions of two ordinal
variables, arXiv:2611.24439 (2026), \doi{10.48550/arXiv.2611.24439}.
\end{thebibliography}

\end{document}
